\documentclass[11pt]{article}

\usepackage[margin=1in]{geometry}
\usepackage{amsmath, amssymb}
\usepackage{xcolor}
\usepackage{listings}
\usepackage{hyperref}
\usepackage{array}
\usepackage{booktabs}

\hypersetup{
    colorlinks=true,
    linkcolor=blue!60!black,
    urlcolor=blue!60!black,
    citecolor=blue!60!black
}

\lstdefinestyle{bugpython}{
  language=Python,
  basicstyle=\ttfamily\small,
  keywordstyle=\color{blue!60!black},
  stringstyle=\color{green!40!black},
  commentstyle=\color{gray!70!black},
  showstringspaces=false,
  breaklines=true,
  frame=single,
  rulecolor=\color{gray!30},
  columns=fullflexible,
  keepspaces=true
}

\title{SymPy Bug Report}
\author{}
\date{}

\begin{document}

\maketitle

\section{Bug 36: Wrong-Sign Antiderivative on the Negative Real Branch}

\subsection{Minimal Reproducer}

\medskip

\begin{lstlisting}[style=bugpython]
from sympy import N, diff, integrate, sqrt, symbols

x = symbols("x")
integrand = 1/(x*sqrt(x**2 + 1))
F = integrate(integrand, x)
dF = diff(F, x)

print("antiderivative:", F)
print("derivative:", dF)
print("integrand at x=-2:", N(integrand.subs(x, -2), 50))
print("derivative at x=-2:", N(dF.subs(x, -2), 50))
print("difference at x=-2:", N((dF - integrand).subs(x, -2), 50))
\end{lstlisting}

\subsection{Observed Behavior}

\medskip

\begin{lstlisting}[style=bugpython]
antiderivative: -asinh(1/x)
derivative: 1/(x**2*sqrt(1 + x**(-2)))
integrand at x=-2: -0.22360679774997896964091736687312762354406183596115
derivative at x=-2: 0.22360679774997896964091736687312762354406183596115
difference at x=-2: 0.44721359549995793928183473374625524708812367192231
\end{lstlisting}

SymPy returns \(-\operatorname{asinh}(1/x)\). At the regular real point \(x=-2\), the derivative of this returned expression is positive, while the original integrand is negative. The mismatch is not a formatting issue: differentiating an indefinite integral should recover the integrand locally on the branch where the returned expression is being used.

\subsection{Expected Behavior}

For
\[
  f(x)=\frac{1}{x\sqrt{x^2+1}},
\]
SymPy should return an antiderivative whose derivative equals \(f(x)\) on the negative real interval as well as on the positive real interval, or it should return a result carrying the necessary branch condition. A single unconditional expression whose derivative has the opposite sign for \(x<0\) is mathematically incorrect.

\subsection{Mathematical Explanation}

SymPy's returned expression is
\[
  F(x)=-\operatorname{asinh}\!\left(\frac{1}{x}\right).
\]
Differentiating gives
\[
  F'(x)
  =
  \frac{1}{x^2\sqrt{1+x^{-2}}}.
\]
For real \(x\ne 0\),
\[
  \sqrt{1+x^{-2}}
  =
  \sqrt{\frac{x^2+1}{x^2}}
  =
  \frac{\sqrt{x^2+1}}{|x|}.
\]
Therefore
\[
  F'(x)
  =
  \frac{|x|}{x^2\sqrt{x^2+1}}
  =
  \frac{1}{|x|\sqrt{x^2+1}}.
\]
On the positive real branch this agrees with \(1/(x\sqrt{x^2+1})\). On the negative real branch it is the negative of the requested integrand:
\[
  \frac{1}{|x|\sqrt{x^2+1}}
  =
  -\frac{1}{x\sqrt{x^2+1}}
  \qquad (x<0).
\]
At \(x=-2\), this gives \(F'(-2)=1/(2\sqrt{5})\), while the integrand is \(-1/(2\sqrt{5})\). The returned expression is thus a one-branch antiderivative, not an equivalent global antiderivative up to an additive constant.

\subsection{Source-Level Diagnosis}

The diagnosis narrows the failure to the indefinite Meijer integration path in
\[
  \text{\nolinkurl{sympy/integrals/meijerint.py}},
\]
specifically \texttt{\_meijerint\_indefinite\_1} around lines 1704--1740, with the closest location identified near line 1736. The default \texttt{integrate} path in \(\text{\nolinkurl{sympy/integrals/integrals.py}}\) tries other methods first; for this integrand, the artifact reports that \texttt{heurisch}, manual integration, and Risch do not produce the final answer, and \texttt{integrate(..., meijerg=False)} remains unevaluated. The result comes from \texttt{meijerint\_indefinite}.

The traced call path is \texttt{integrate} to \texttt{meijerint\_indefinite} to \texttt{\_meijerint\_indefinite\_1} to \texttt{\_rewrite1}. For the reproducer, \texttt{\_rewrite1} represents the integrand as an \(x^{-1}\) power factor times a Meijer \(G\)-function of \(x^2\):
\[
  (1,\;1/x,\;[(1/\sqrt{\pi},\,0,\,
  \operatorname{meijerg}(((1/2),()),((0),()),x^2))],\;\mathrm{True}).
\]
The problematic behavior appears in the branch-blind substitution and reconstruction step:

\begin{lstlisting}[style=bugpython]
a, b = _get_coeff_exp(g.argument, x)
_, c = _get_coeff_exp(po, x)
...
fac_ = fac * C * x**(1 + c) / b
rho = (c + 1)/b
...
r = hyperexpand(r.subs(t, a*x**b), place=place)
res += powdenest(fac_*r, polar=True)
\end{lstlisting}

For this input, the substitution is effectively \(t=x^2\) while the prefactor contains \(1/x\). That substitution identifies the positive and negative real branches of \(x\). The resulting Meijer antiderivative is then hyperexpanded to \(-\operatorname{asinh}(1/x)\) without a condition recording that the expression differentiates correctly only on one real branch. This explains the observed mathematical failure: the square-root simplification hidden in the derivative uses \(|x|\), while the original integrand uses \(x\).

The suggested fix direction in the diagnosis is to avoid branch-losing substitutions such as \(x\mapsto x^2\) in the indefinite Meijer path when odd powers such as \(1/x\) are present, unless the method can return a piecewise or otherwise branch-conditioned antiderivative. At minimum, the Meijer result should be checked against the original integrand under real branch assumptions, or declined when the derivative check is branch-dependent.

\subsection{Additional Failing Instances}

The same sign error occurs for the family
\[
  \int \frac{dx}{x\sqrt{x^2+a}},
  \qquad a\in\{1,2,3,4,5,9\}.
\]
For positive \(a\), SymPy returns the corresponding one-branch expression
\[
  F_a(x)=-\frac{\operatorname{asinh}(\sqrt{a}/x)}{\sqrt{a}},
\]
with the obvious simplifications when \(a\) is a square. Differentiation gives the positive value \(1/(2\sqrt{4+a})\) at \(x=-2\), while the original integrand has value \(-1/(2\sqrt{4+a})\). The expected behavior is the same throughout the family because \(\sqrt{x^2+a}/x\) and \(\sqrt{1+a/x^2}\) differ by the sign of \(x\) on the real line.

\begin{center}
\renewcommand{\arraystretch}{1.3}
\begin{tabular}{@{}>{\raggedright\arraybackslash}p{0.44\linewidth}>{\raggedright\arraybackslash}p{0.23\linewidth}>{\raggedright\arraybackslash}p{0.23\linewidth}@{}}
\toprule
\textbf{Input / Expression} & \textbf{SymPy behavior} & \textbf{Expected behavior} \\
\midrule
\(a=1:\;1/(x\sqrt{x^2+1})\) & \(F'(-2)=1/(2\sqrt{5})\) & \(-1/(2\sqrt{5})\) \\
\(a=2:\;1/(x\sqrt{x^2+2})\) & \(F'(-2)=1/(2\sqrt{6})\) & \(-1/(2\sqrt{6})\) \\
\(a=3:\;1/(x\sqrt{x^2+3})\) & \(F'(-2)=1/(2\sqrt{7})\) & \(-1/(2\sqrt{7})\) \\
\(a=4:\;1/(x\sqrt{x^2+4})\) & \(F'(-2)=1/(4\sqrt{2})\) & \(-1/(4\sqrt{2})\) \\
\(a=5:\;1/(x\sqrt{x^2+5})\) & \(F'(-2)=1/6\) & \(-1/6\) \\
\(a=9:\;1/(x\sqrt{x^2+9})\) & \(F'(-2)=1/(2\sqrt{13})\) & \(-1/(2\sqrt{13})\) \\
\bottomrule
\end{tabular}
\end{center}

\subsection{Related Issues Assessment}

The best-effort deduplication check found documentation and background material related to the general family of branch-sensitive substitutions, including SymPy documentation warning that quadratic substitutions are acceptable only when the transformed integrand does not depend on the sign of the solutions. It did not identify a public SymPy issue, pull request, Stack Overflow report, mailing-list post, or release note for this exact integrand, the output \(-\operatorname{asinh}(1/x)\), or the derivative sign mismatch at a negative real point. The deduplication verdict is therefore a new instance of a known failure mode: the broader branch-loss mechanism is understood, but this exact reproducible indefinite-integration case appears previously unreported and remains individually actionable as a focused issue and regression test.

\subsection{Proposed SymPy Regression Test}

\medskip

\begin{lstlisting}[style=bugpython]
import pytest

from sympy import N, diff, integrate, sqrt, symbols


@pytest.mark.parametrize("a", [1, 2, 3, 4, 5, 9])
def test_integrate_x_sqrt_x2_plus_a_negative_branch_sign(a):
    x = symbols("x")
    integrand = 1/(x*sqrt(x**2 + a))
    F = integrate(integrand, x)
    difference = N((diff(F, x) - integrand).subs(x, -2), 50)

    assert abs(complex(difference)) < 1e-40
\end{lstlisting}

\end{document}
